% Section I: Introduction
\section{Introduction}
\label{sec:introduccion}
\IEEEPARstart{V}{ision} Transformers (ViTs) \cite{dosovitskiy2021image} reach the state of the art in image classification when pretrained at large scale, but their compact variants, trained from scratch on limited data, generalize worse than convolutional networks. CNNs build in inductive biases such as spatial locality and translation equivariance; a ViT splits the image into linearly projected patches and relies on the network to learn the spatial topology from position encodings. With enough data, attention learns local patterns in the first layers \cite{raghu2021vision}. With little data and a small model, such as the 119,290-parameter model of this work, it does not: the mean attention distance is 14 to 17 pixels on 32-pixel images in every layer.

Compact architectures are of interest when both data and compute are scarce: they must learn from small datasets, from scratch and at a low training cost. This work studies that regime (data efficiency, convergence speed and calibration). Deployment on edge devices (\emph{Edge AI}) is a further motivation, but here we only measure the inference cost on a desktop CPU; embedded hardware is not evaluated (Section~\ref{sec:res-costos}). Heterogeneous attention mechanisms (HAM), which combine more than one way of mixing the values within a head, have been proposed to improve the attention of these models. In previous work \cite{martinez2025reconfiguration} we proposed two such mechanisms for the same ViT: HAM A and HAM B, which combine in each head the output of self-attention with that of a pooling-type additive attention (\emph{soft attention}), through the Hadamard product and the element-wise sum, respectively. HAM B reached 81.07~\% accuracy on CIFAR-10, 1.3 points above the reference ViT (79.77~\%), and peaked earlier. That evaluation used a single run per variant, a constant learning rate of 0.01 and the best of 400 epochs measured on the test set; at the last epoch HAM B was below the reference (76.93 versus 79.02~\%). This work revisits those results with a controlled protocol and separates two ways of understanding heterogeneity: combining two similarity functions over the same queries and keys, or adding to the attention output terms with an inductive bias that attention lacks. The most complete configuration, which we call the \emph{triad}, brings three biases together in each head: global pairwise similarity, a global summary of the image and spatial neighborhood. The central question is how much of its gain is due to each bias and how much to the parameters it adds.

The term that turns out to matter is not new: the depthwise convolution over the values is the locally enhanced positional encoding (LePE) of CSWin \cite{dong2022cswin}. The contribution of this work is therefore not a new architecture but a controlled measurement, with three parts: a well-controlled negative result on fusing similarity functions, the contribution of each inductive bias separated from the parameters it adds, and methodological evidence on how many seeds are needed to compare attention variants in this regime.

The contributions are the following:
\begin{itemize}
  \item A controlled evaluation of two families of heterogeneous mechanisms in a small ViT: the same architecture, paired seeds with deterministic training, 16 seeds in the main comparison, equal-capacity controls and correction for multiple comparisons, on CIFAR-10, CIFAR-100, SVHN and PlantVillage, and on Food-101 for the fusions.
  \item A well-controlled negative result for the original HAM hypothesis, and its explanation: mixing scores before the \emph{softmax} multiplies the distributions, so one branch decides, and no fusion separates from a control with the same parameters.
  \item A measurement of the contribution of a local term, a depthwise convolution over the values of each head: $+5.31$ points on CIFAR-10 over an equal-capacity control ($+5.64$ over the base model), with 4.2~\% more parameters (3.1~\% of them effective) and 3.4~\% more MACs, better calibration and robustness, and 3.7 times less GPU time to reach 79~\% validation accuracy, close to the final value of the base model.
  \item A decomposition of the triad with equal-capacity controls and controls for the fusion projection: the neighborhood explains almost all of the gain; the global summary branch adds $+0.61$ points over an equal-capacity control in a contrast declared in advance and extended to 48 seeds ($+0.67$ on the 32 new seeds), and improves calibration. The triad, however, does not meet the cost criterion set in advance.
  \item A test of the scope of the local term: it holds on four datasets, with CutMix (which does not induce local heads at this scale) and from 119,000 to 3.6 million parameters, and it is compared with four published locality mechanisms for ViTs, reimplemented at the same budget, and with a ResNet of the same size.
  \item An analysis of its effects on attention and methodological evidence for this regime: the local term gives rise to two training modes, which inflate the variance of the paired differences, so half-point effects need tens of paired seeds to be resolved.
\end{itemize}

To consider that a variant works, four criteria were set in advance: a significant gain over an equal-capacity control, a balanced contribution of each branch, the same sign on another dataset and a cost of at most 1.5 times the epoch time of the base model (Section~\ref{sec:analisis}). No variant meets all four; the triad fails the cost criterion (Section~\ref{sec:res-criterios}). The rest of the article is organized as follows: Section~\ref{sec:trabajos} reviews related work; Sections~\ref{sec:metodologia} and~\ref{sec:diseno} describe the mechanisms and the experimental design; Section~\ref{sec:resultados} presents the results; Section~\ref{sec:discusion} discusses them, and Section~\ref{sec:conclusiones} concludes. The ablations of the output terms, the attention maps and a sensitivity analysis of the threshold that separates the two training modes are given as Supplementary Material.
